%% 5-13 積分定数は「はじめの状態」（第11回：ページ制限なし。C がなぜ「はじめの位置」なのかを、表・グラフ・式の3つで見せる）
\bkchapter{5}{13}{積分定数は「はじめの状態」}{\bklevel{3}}{積分定数は「はじめの状態」}
\bkmichishirube{5-5,5-9,5-11}{6-1,6-3}

\begin{goal}
\begin{itemize}[label={\color{bkTheme}●}, leftmargin=1.6\zw, itemsep=0pt, topsep=1pt]
\item \bklevel{3}\enspace 積分定数 $C$ を、はじめの状態（初期条件）から決められる。
\item \bklevel{3}\enspace 加速度から速度、速度から位置を、積分で求められる。
\item \bklevel{3}\enspace $x=x_0+v_0t+\dfrac{1}{2}at^2$ を、積分で導ける。
\end{itemize}
\end{goal}

\begin{hatena}
\noindent\begin{minipage}[c]{0.58\linewidth}\raggedright
AさんとBさんは、同じ時刻に、同じ速さで、同じ向きに歩き出しました。1時間のあいだ、2人の速さはいつも同じでした。

では、1時間後、2人は同じ場所にいるでしょうか。
\end{minipage}\hfill
\begin{minipage}[c]{0.38\linewidth}
\begin{tikzpicture}[x=1mm, y=1mm]
  \draw[bkGray, line width=1pt] (0,0) -- (40,0);
  \draw[bkGray, line width=1pt] (0,12) -- (40,12);
  \bkperson{(2,0.5)}{10}{bkRoute2}
  \bkperson{(14,12.5)}{10}{bkRoute4}
  \draw[bk arrow, bkGray, line width=0.9pt] (8,6) -- (18,6);
  \draw[bk arrow, bkGray, line width=0.9pt] (20,18) -- (30,18);
  \node[font=\sffamily\scriptsize, anchor=east] at (-1.5,5) {A};
  \node[font=\sffamily\scriptsize, anchor=east] at (10.5,17) {B};
\end{tikzpicture}
\end{minipage}
\end{hatena}

速さがいつも同じなら、1時間に進む距離は同じです。でも、Bさんがはじめから\qty{1}{km}先にいたなら、1時間後も\qty{1}{km}先にいます。速さだけでは、位置は決まりません。「はじめにどこにいたか」が必要です。この章では、5-11で出てきた積分の $+C$ が、この「はじめの状態」を表していることを確かめます。

\clearpage
%% ---------------- やってみよう・思い出そう ----------------
\begin{yattemiyou}[同じ速度、ちがうスタート]
速度が $v=2t$（m/s）の物体が3つあります。スタートの位置は、Aが\qty{0}{m}、Bが\qty{3}{m}、Cが\qty{-2}{m}です（Cはスタートの線より\qty{2}{m}後ろ）。
\begin{enumerate}[label={\bkcircled[bkGreen]{\arabic*}}, leftmargin=2.4\zw, itemsep=1pt, topsep=3pt]
\item 0秒から $t$ 秒までに進む距離を、$v$-$t$ グラフの面積で求めよう。
\item それぞれの位置（スタートの位置＋進んだ距離）を、$t=0$〜3 で表にしよう。
\item 横軸を $t$、たて軸を位置にして、3つのグラフをかこう（$x$-$t$ グラフ）。
\end{enumerate}
\begin{bkcenter}
\begin{bktable}{colspec={Q[l,m,wd=34mm] *{4}{Q[c,m,wd=12mm]}}, rowsep=1.5pt}
$t$（s） & 0 & 1 & 2 & 3 \\
進んだ距離 $t^2$（m） & 0 & 1 & 4 & 9 \\
Aの位置（m） & 0 & 1 & 4 & 9 \\
Bの位置（m） & 3 & 4 & 7 & 12 \\
Cの位置（m） & $-2$ & $-1$ & 2 & 7 \\
\end{bktable}
\end{bkcenter}
\end{yattemiyou}

\noindent ①の面積は、底辺 $t$、高さ $2t$（$t$ 秒の速度）の三角形なので、$\dfrac{1}{2}\times t\times2t=t^2$ です（5-9）。5-11の積分で書くと、
\[ \int_0^t 2s\,\dd s=\Bigl[\,s^2\,\Bigr]_0^t=t^2-0=t^2 \]
です（$s$ は0から $t$ まで動く時刻。上端の $t$ と区別するため、別の文字にしています）。どちらの方法でも、進んだ距離は $t^2$ です。

3つのグラフは、同じ形の放物線を、上下にずらしたものになります。どの時刻でも、3つのグラフの傾き（速度）は同じです。それでも位置がちがうのは、スタートの位置がちがうからです。

\begin{omoidasou}
\bkrecall{5-11}{微分すると $f(x)$ になる関数が原始関数。$\displaystyle\int f(x)\,\dd x=F(x)+C$。定数は微分すると0なので、原始関数には $+C$ の自由がある。}
\bkrecall{5-5}{速度は位置を時間で微分したもの、加速度は速度を時間で微分したもの。落ちる物の加速度は $g=\qty{9.8}{m/s^2}$。}
\bkrecall{5-9}{$v$-$t$ グラフの面積は、位置の変化（変位）。}
\bkrecall{2-2}{マイナスは「反対の向き」を表す。$+$ の向きは、はじめに決めて変えない。}
\bkrecall{3-10}{$t^2+2t-15=0$ は、かけて $-15$、たして2になる2つの数（5と $-3$）を見つけて、$(t+5)(t-3)=0$ と因数分解する。かけて0になるのは、どちらかが0のときなので、$t=-5$ か $t=3$。}
\end{omoidasou}

%% ---------------- 数学の道具 ----------------
\bksection{doguu}{数学の道具}

\bksubsection{1}{速度を積分すると、位置が「$+C$」つきでわかる}

\label{bk:5-13:C}位置 $x(t)$ を微分すると速度 $v(t)$ です（5-5）。逆に言うと、位置 $x(t)$ は「微分すると $v(t)$ になる関数」、つまり $v(t)$ の原始関数です（5-11）。だから
\[ x(t)=\int v(t)\,\dd t=F(t)+C \]
と書けます。$v=2t$ なら、$t^2$ を微分すると $2t$ なので、$x=t^2+C$ です。

\noindent\textbf{なぜ $+C$ がつくの？}\quad $t^2+3$ も $t^2-2$ も、微分すると $2t$ です。定数（3や $-2$）は時間がたっても変わらないので、変わる速さは0だからです。表で確かめます。

\begin{bkcenter}
\begin{bktable}{colspec={Q[l,m,wd=40mm] *{4}{Q[c,m,wd=9mm]} Q[c,m,wd=26mm]}, rowsep=1.5pt, row{1}={bg=bkTheme!10}}
$t$（s） & 0 & 1 & 2 & 3 & 1秒ごとのふえ方 \\
$x=t^2$（$C=0$） & 0 & 1 & 4 & 9 & 1，3，5 \\
$x=t^2+3$（$C=3$） & 3 & 4 & 7 & 12 & 1，3，5 \\
$x=t^2-2$（$C=-2$） & $-2$ & $-1$ & 2 & 7 & 1，3，5 \\
\end{bktable}
\bktabcaption{$C$ がちがっても、1秒ごとのふえ方（動き方）はまったく同じ。}
\end{bkcenter}

$C$ の値によって、グラフは上下にずれるだけです。どの $C$ でも、速度（傾き）は同じ $2t$ です。だから、速度だけでは $C$ は決まりません。

\noindent\begin{minipage}[c]{0.5\linewidth}
\begin{tikzpicture}
\begin{axis}[bkplot, width=62mm, height=48mm, xmin=0, xmax=3.4, ymin=-2.6, ymax=13, xtick={1,2,3}, xticklabels={}, ytick={-2,3,9}, xlabel={$t$（s）}, ylabel={$x$（m）}]
  \addplot[bkcurve, domain=0:2.75] {x^2};
  \addplot[bkcurve2, domain=0:2.75] {x^2+3};
  \addplot[bkcurve3, domain=0:2.75] {x^2-2};
  % 目もりの数字は、緑の線（t=1 で −1）と重ならないよう、−2 の線より下に置く
  \node[font=\footnotesize, text=bkInk, anchor=north] at (axis cs:1,-3.2) {$1$};
  \node[font=\footnotesize, text=bkInk, anchor=north] at (axis cs:2,-3.2) {$2$};
  \node[font=\footnotesize, text=bkInk, anchor=north] at (axis cs:3,-3.2) {$3$};
  \node[font=\sffamily\scriptsize, anchor=west, text=bkOrange!80!black, fill=white, inner sep=1pt] at (axis cs:2.85,10.6) {$C=3$（B）};
  \node[font=\sffamily\scriptsize, anchor=west, text=bkBlue, fill=white, inner sep=1pt] at (axis cs:2.85,7.6) {$C=0$（A）};
  \node[font=\sffamily\scriptsize, anchor=west, text=bkGreen!70!black, fill=white, inner sep=1pt] at (axis cs:2.85,5.6) {$C=-2$（C）};
\end{axis}
\end{tikzpicture}
\end{minipage}\hfill
\begin{minipage}[c]{0.46\linewidth}
$x=t^2+C$ に $t=0$ を入れると、
\[ x(0)=0^2+C=C \]
です。つまり、\textbf{積分定数 $C$ は、はじめ（$t=0$）の位置}を表しています。

はじめの位置 $x_0$ がわかれば、$C=x_0$ と決まり、どの時刻の位置も1つに決まります。
\end{minipage}

\noindent\textbf{\bkhatenaref の答え}\quad AさんとBさんは、速さがいつも同じなので、同じ時間に同じ距離だけ進みます。だから、1時間後の2人の差は、はじめの差のままです。Bさんがはじめ\qty{1}{km}先にいたなら1時間後も\qty{1}{km}先、同じ場所から出発したなら同じ場所にいます。位置の式で言えば、2人の式は $C$（はじめの位置）だけがちがうのです。

\bksubsection{2}{はじめの状態：初期条件}

\label{bk:5-13:shoki}$C$ を決めるのに使う「ある時刻の値」を、\textbf{初期条件}（しょきじょうけん）といいます。たいていは $t=0$ の値ですが、ほかの時刻の値でも $C$ は決められます。\bkidx{しょきじょうけん}{初期条件}

\noindent\begin{minipage}[c]{0.5\linewidth}
\begin{tikzpicture}
\begin{axis}[bkplot, width=62mm, height=48mm, xmin=0, xmax=3.4, ymin=-1, ymax=16, xtick={1,2}, ytick={4,8,12}, xlabel={$t$（s）}, ylabel={$x$（m）}]
  \foreach \c in {0,6} {\edef\tmp{\noexpand\addplot[bkthin, domain=0:2.75] {x^2+\c};}\tmp}
  \addplot[bkcurve2, domain=0:2.75] {x^2+3};
  \addplot[only marks, mark=*, mark size=2pt, bkInk] coordinates {(1,4)};
  \node[font=\sffamily\scriptsize, anchor=west, text=bkOrange!80!black, fill=white, inner sep=1pt] at (axis cs:2.82,10.6) {$C=3$};
\end{axis}
\end{tikzpicture}
\end{minipage}\hfill
\begin{minipage}[c]{0.46\linewidth}\raggedright
たとえば、速度 $v=2t$ の物体が、「1秒のとき\qty{4}{m}の所にいた」とわかったとします。$x=t^2+C$ に $t=1$、$x=4$ を入れると
\[ 4=1^2+C,\qquad C=4-1=3 \]
グラフでは、上下にずれた放物線のなかから、点 $(1,\,4)$（黒い点）を通る1本（オレンジ）を選んだことになります。
\end{minipage}

\noindent\textbf{$C$ を決める4つの手順}\quad どの問題でも、次の順に進めば $C$ が決まります。右の列は、いまの例（$v=2t$、1秒のとき\qty{4}{m}）です。
\begin{bkcenter}
\begin{bktable}{colspec={Q[c,m,wd=8mm] Q[l,m,wd=60mm] Q[l,m,wd=50mm]}, rowsep=2pt, row{1}={bg=bkTheme!10}}
 & すること & いまの例 \\
① & 積分して、$+C$ をつける & $x=t^2+C$ \\
② & 初期条件の時刻と値を入れる & $t=1$、$x=4$ を入れて $4=1^2+C$ \\
③ & $C$ について解く & 両辺から1をひいて $C=3$ \\
④ & 確かめる（微分してもとの式、初期条件を入れて合うか） & $(t^2+3)'=2t$、$1^2+3=4$ \\
\end{bktable}
\bktabcaption{①で $+C$ を忘れると、②③ができなくなる。}
\end{bkcenter}

このあと、例題1は $t=0$ の値で、例題2は $t=1$ の値で $C$ を決めます。

\bksubsection{3}{加速度から位置まで：2回積分する}

\label{bk:5-13:toukasoku}加速度 $a$ が一定のとき、2回積分して位置を求めます。5-5で予告した式 $x=x_0+v_0t+\dfrac{1}{2}at^2$ を、ここで導きましょう。

\noindent\textbf{1回め：加速度から速度}\quad 速度を微分すると加速度なので、速度は加速度の原始関数です。$a$ は定数で、$at$ を微分すると $a$ になる（5-4）ので、
\[ v=\int a\,\dd t=at+C_1 \]
$t=0$ を入れると $v(0)=a\times0+C_1=C_1$ なので、$C_1$ は\textbf{はじめの速度}です。これを $v_0$（ブイ・ゼロ）と書くと、$C_1=v_0$ で
\[ v=v_0+at \]
「はじめの速度に、1秒ごとに $a$ ずつたしていく」という意味です。

\noindent\textbf{2回め：速度から位置}\quad 位置は、速度の原始関数です。$v_0$ も $a$ も定数なので、項ごとに積分します（5-11）。$v_0$ を積分すると $v_0t$、$at$ を積分すると $a\times\dfrac{t^2}{2}$（指数を1ふやして、その数でわる）です。
\[ x=\int(v_0+at)\,\dd t=v_0t+\frac{1}{2}at^2+C_2 \]
$t=0$ を入れると $x(0)=C_2$ なので、$C_2$ は\textbf{はじめの位置} $x_0$ です。だから
\[ x=x_0+v_0t+\frac{1}{2}at^2 \]

\noindent\textbf{数と図で確かめる}\quad $a=\qty{2}{m/s^2}$、$v_0=\qty{3}{m/s}$、$x_0=\qty{1}{m}$ の場合を表にします。式は $v=3+2t$、$x=1+3t+t^2$ です。

\begin{bkcenter}
\begin{bktable}{colspec={Q[l,m,wd=50mm] *{4}{Q[c,m,wd=11mm]}}, rowsep=1.5pt}
$t$（s） & 0 & 1 & 2 & 3 \\
速度 $v=3+2t$（m/s） & 3 & 5 & 7 & 9 \\
1秒ごとの台形の面積（m） & & 4 & 6 & 8 \\
位置 $x=1+3t+t^2$（m） & 1 & 5 & 11 & 19 \\
\end{bktable}
\bktabcaption{位置は、はじめの1 m に、$v$-$t$ グラフの台形の面積をつみ上げたもの。}
\end{bkcenter}

\noindent\begin{minipage}[c]{0.46\linewidth}
\begin{tikzpicture}
\begin{axis}[bkplot, width=56mm, height=44mm, xmin=0, xmax=3.6, ymin=0, ymax=11, xtick={1,2,3}, ytick={3,5,7,9}, xlabel={$t$（s）}, ylabel={$v$（m/s）}]
  \addplot[fill=bkBlue!18, draw=bkBlue!60, line width=0.4pt] coordinates {(0,0) (0,3) (1,5) (1,0)} -- cycle;
  \addplot[fill=bkGreen!22, draw=bkGreen!60!black, line width=0.4pt] coordinates {(1,0) (1,5) (2,7) (2,0)} -- cycle;
  \addplot[fill=bkOrange!25, draw=bkOrange!70, line width=0.4pt] coordinates {(2,0) (2,7) (3,9) (3,0)} -- cycle;
  \addplot[bkcurve, domain=0:3.4] {3+2*x};
  \node[font=\sffamily\small] at (axis cs:0.5,1.5) {4};
  \node[font=\sffamily\small] at (axis cs:1.5,2.2) {6};
  \node[font=\sffamily\small] at (axis cs:2.5,3) {8};
\end{axis}
\end{tikzpicture}
\end{minipage}\hfill
\begin{minipage}[c]{0.52\linewidth}\raggedright
速度は1秒ごとに2ずつふえます（加速度2）。0〜1秒の $v$-$t$ グラフは上底3、下底5の台形（青）で、面積は $(3+5)\times1\div2=4$。はじめの位置1に4をたすと5で、式の $x(1)=1+3+1=5$ と合います。次の1秒は緑の台形6をたして11、その次はオレンジの8をたして19です。
\end{minipage}

\noindent $\dfrac{1}{2}at^2$ の $\dfrac{1}{2}$ は、5-9の三角形の面積の「÷2」と同じものです。台形を「長方形 $v_0\times t$」と「三角形 $\dfrac{1}{2}\times t\times at$」に分けると、$v_0t+\dfrac{1}{2}at^2$ になるからです。

2回積分したので、積分定数が2つ出てきました。それを決めるのに、はじめの位置 $x_0$ とはじめの速度 $v_0$ の、2つの初期条件が必要です。たとえば同じ $a=2$ でも、$v_0$ が3と0では、1秒後の速度が5と2でちがいます。$v_0$ だけわかっても、$x_0$ がわからなければ、どこにいるかは決まりません。

\bksubsection{4}{落ちる物・投げ上げた物}

\label{bk:5-13:rakka}\noindent\begin{minipage}[c]{0.6\linewidth}\raggedright
上向きを＋と決めます（2-2）。位置 $x$ は、地面からの高さです。重力は、物を1秒ごとに下向きへ\qty{9.8}{m/s}ずつ速くします。下向きは「＋の反対の向き」なので、加速度は $a=-9.8$ です。

高さ $x_0$ から、上向きの速さ $v_0$ で投げると、③の式に $a=-9.8$ を入れて
\[ v=v_0-9.8t,\qquad x=x_0+v_0t-4.9t^2 \]
4.9は、$\dfrac{1}{2}at^2$ の $\dfrac{1}{2}\times9.8=4.9$ から来ています。
\end{minipage}\hfill
\begin{minipage}[c]{0.36\linewidth}\centering
\begin{tikzpicture}[x=1mm, y=1mm]
  \fill[bkGray!35] (-2,-3) rectangle (40,0);
  \draw[bkInk, line width=0.7pt] (-2,0) -- (40,0);
  \node[font=\sffamily\scriptsize, anchor=north] at (24,-4) {地面（$x=0$）};
  \draw[bk arrow, bkInk, line width=0.8pt] (4,0) -- (4,40);
  \node[font=\sffamily\scriptsize, anchor=south] at (4,40.5) {＋（上）};
  \draw[bkInk, line width=0.5pt] (2.5,28) -- (5.5,28);
  \node[font=\sffamily\scriptsize, anchor=east] at (2,28) {$x_0$};
  \fill[bkOrange!70] (14,28) circle (2.6);
  \draw[bkOrange!80!black, line width=0.5pt] (14,28) circle (2.6);
  \draw[bk arrow, bkBlue, line width=0.9pt] (14,31.5) -- (14,39);
  \node[font=\sffamily\scriptsize, text=bkBlue, anchor=west] at (15.5,36) {$v_0$（上向き）};
  \draw[bk arrow, bkRed, line width=0.9pt] (14,24.5) -- (14,13);
  \node[font=\sffamily\scriptsize, text=bkRed, anchor=west, align=left] at (15.5,18) {加速度\\$a=-9.8$};
\end{tikzpicture}
\end{minipage}

\par\medskip
\noindent\textbf{例}\quad 高さ\qty{20}{m}から、ボールを静かに落とします。「静かに落とす」は、はじめの速度が0ということです。$x_0=20$、$v_0=0$ なので、位置は $x=20-4.9t^2$ です。地面は高さ0なので、$x=0$ になる時刻を求めます。
\begin{align*}
20-4.9t^2&=0\\
4.9t^2&=20\qquad\text{（両辺に $4.9t^2$ をたす）}\\
t^2&=20\div4.9\approx4.08\qquad\text{（両辺を4.9でわる）}
\end{align*}
2回かけて4.08になる正の数は、電卓の√キーで $\sqrt{4.08}\approx2.02$ です（2-8）。$t$ は時間なので、正のほうだけを使います。地面に着くのは約\qty{2.0}{s}後です。

\bksubsection{5}{定積分で書くと「はじめ＋たまった変化」}

\label{bk:5-13:teisekibun}位置の変化は、$v$-$t$ グラフの面積でした（5-9）。はじめの位置にそれをたせば、いまの位置です。
\[ x(t)=x_0+\int_0^t v(s)\,\dd s \]
（$s$ は0から $t$ まで動く時刻です。）おこづかいでたとえると、「いまの貯金＝はじめの貯金＋その後にたまった分」です。たまった分がわかっても、はじめの貯金がわからなければ、いまの貯金はわかりません。$C$ や $x_0$ が「はじめの貯金」にあたります。

「いまの量＝はじめの量＋たまった変化」は、水そうの水の量、温度、電池の電気の量など、変化の速さがわかる量すべてに使えます（5-12で求めた「たまった量」に、はじめの量をたせばよい）。

\begin{bkfigure}
\begin{tikzpicture}[x=1mm, y=1mm, m/.style={draw=bkTheme, line width=1pt, fill=white, rounded corners=2mm, minimum width=28mm, minimum height=11mm, align=center, font=\sffamily\small}]
  \node[m] (a) at (0,0) {加速度 $a$};
  \node[m] (v) at (48,0) {速度 $v$};
  \node[m] (x) at (96,0) {位置 $x$};
  \draw[bk arrow, bkBlue] (a) -- node[above, font=\sffamily\scriptsize] {積分} node[below, font=\sffamily\scriptsize] {$+v_0$} (v);
  \draw[bk arrow, bkBlue] (v) -- node[above, font=\sffamily\scriptsize] {積分} node[below, font=\sffamily\scriptsize] {$+x_0$} (x);
\end{tikzpicture}
\bkfigcaption{積分するたびに、はじめの状態を1つ使う。}
\end{bkfigure}

\bksubsection{6}{表計算で、はじめの状態から先を計算する}

\label{bk:5-13:hyou}式で積分できなくても、5-8の方法で、はじめの状態から先を計算できます。「短い時間 $\Delta t$ のあいだは、速度も加速度もほとんど変わらない」と考えて、次のようにくり返します。
\[ \text{次の速度}=\text{今の速度}+a\times\Delta t,\qquad \text{次の位置}=\text{今の位置}+\text{今の速度}\times\Delta t \]
投げ上げたボール（$a=-10$、$v_0=20$、$x_0=0$）を $\Delta t=0.5$ で計算します。たとえば3行めは、速度 $20+(-10)\times0.5=15$、位置 $0+20\times0.5=10$ です。
\begin{bkcenter}
\bksheet{5}{20}{%
1 & $t$ & 速度 $v$ & 位置 $x$（予想） & $20t-5t^2$ & ずれ \\
2 & 0 & 20 & 0 & 0 & 0 \\
3 & 0.5 & 15 & 10 & 8.75 & 1.25 \\
4 & 1.0 & 10 & 17.5 & 15 & 2.5 \\
5 & 1.5 & 5 & 22.5 & 18.75 & 3.75 \\
6 & 2.0 & 0 & 25 & 20 & 5 \\}
\bktabcaption{はじめの2つの値（$v_0=20$、$x_0=0$）がないと、計算を始められない。}
\end{bkcenter}
表の2行めの「はじめの速度」と「はじめの位置」が、積分定数 $C_1$、$C_2$ にあたります。予想の位置が少し大きいのは、0.5秒のあいだにへっていく速度を、はじめの（大きい）速度のままとして計算したからです。ずれは $\Delta t$ を小さくするとへります（5-8）。

\begin{column}{ゲームの中で、ボールはどう動いている？}
ボールを投げるゲームの中では、コンピュータが⑥と同じ計算をしています。画面は1秒に60回かきかえられるので、$\Delta t$ は $\dfrac{1}{60}$ 秒（約0.017秒）です。1回かきかえるたびに「次の速度＝今の速度＋加速度×$\dfrac{1}{60}$」「次の位置＝今の位置＋今の速度×$\dfrac{1}{60}$」を計算します。ボールを投げた瞬間に決まる「投げた場所」と「投げた速度」が、はじめの状態（$x_0$ と $v_0$）です。同じ場所・同じ速度で投げれば、何回やっても同じ所に落ちます。
\end{column}

\begin{column}{未来は、はじめの状態で決まる？}
200年ほど前、フランスのラプラスは「すべての物の位置と速度がわかれば、運動の法則で未来はすべて計算できる」と考えました。この章の考えを、宇宙全体に広げたものです。でも、はじめのわずかなずれが大きく育つ現象（天気など）もあり、未来の予想には限界があります。
\end{column}

\begin{taisetsu}
積分定数 $C$ は、はじめの状態（初期条件）で決まる。\\
加速度が一定なら\quad $v=v_0+at$、\quad $x=x_0+v_0t+\dfrac{1}{2}at^2$\\
いまの量 ＝ はじめの量 ＋ 変化の速さの積分
\end{taisetsu}

\begin{kotoba}
\small
\bkword{初期条件}{積分定数を決めるのに使う、ある時刻（ふつうは $t=0$）の値。}
\bkword{$x_0$、$v_0$}{はじめの位置、はじめの速度。「エックス・ゼロ」「ブイ・ゼロ」と読む。}
\end{kotoba}

%% ---------------- 例題 ----------------
\bksection{reidai}{例題でたしかめよう}

\begin{reidai}[初期条件で $C$ を決める]
速度が $v=6t+2$（m/s）で、$t=0$ の位置が\qty{5}{m}の物体の、位置 $x$ を $t$ の式で表しましょう。\par
\kangaekata 位置は速度の原始関数（①）。項ごとに積分して $+C$ をつけ、$t=0$ の位置で $C$ を決めます。\par
\shiki $6t$ を積分すると $6\times\dfrac{t^2}{2}=3t^2$、2を積分すると $2t$\par
\hspace*{2.6\zw}$\displaystyle x=\int(6t+2)\,\dd t=3t^2+2t+C$\par
\hspace*{2.6\zw}$t=0$ のとき $x=5$ なので、$5=3\times0^2+2\times0+C$ より $C=5$\par
\hspace*{2.6\zw}確かめ：$x=3t^2+2t+5$ を微分すると $6t+2$\par
\kotae $x=3t^2+2t+5$
\end{reidai}

\begin{jibun}
速度 $v=4t-3$（m/s）、$t=0$ で位置\qty{2}{m}の物体の位置を求めましょう。
\bkjans{$x=2t^2-3t+C$、$C=2$。$x=2t^2-3t+2$}
\end{jibun}

\noindent\begin{minipage}[c]{0.48\linewidth}
\begin{tikzpicture}
\begin{axis}[bkplot, width=58mm, height=44mm, xmin=0, xmax=2.4, ymin=0, ymax=25, xtick={1,2}, ytick={5,10,20}, xlabel={$t$（s）}, ylabel={$x$（m）}]
  \addplot[bkcurve, domain=0:2.15] {3*x^2+2*x+5};
  \addplot[only marks, mark=*, mark size=2pt, bkInk] coordinates {(0,5)};
  \node[font=\sffamily\scriptsize, anchor=west, fill=white, inner sep=1pt] at (axis cs:0.06,15) {出発点 $(0,\,5)$};
  \draw[bkGray, line width=0.5pt] (axis cs:0.12,13.6) -- (axis cs:0.03,6);
\end{axis}
\end{tikzpicture}
\end{minipage}\hfill
\begin{minipage}[c]{0.5\linewidth}\raggedright
\bkreidairef{1}の答え $x=3t^2+2t+5$ のグラフです。$t=0$ のとき $x=5$ で、グラフは高さ5から始まります。$C=5$ は、グラフの出発点の高さです。$C$ を変えると、同じ形のまま、出発点だけが上下に動きます。
\end{minipage}

\begin{reidai}[$t=0$ でない初期条件]
速度 $v=3t^2-4$（m/s）、$t=1$ で位置0の物体の位置を求めましょう。\par
\kangaekata 積分して $+C$ をつけ、$t=1$、$x=0$ を入れて $C$ を決めます（②）。\par
\shiki $\displaystyle x=\int(3t^2-4)\,\dd t=3\times\frac{t^3}{3}-4t+C=t^3-4t+C$\par
\hspace*{2.6\zw}$t=1$ のとき $x=0$ なので、$0=1^3-4\times1+C=-3+C$\par
\hspace*{2.6\zw}両辺に3をたして $C=3$\par
\kotae $x=t^3-4t+3$（はじめの位置は $x(0)=\qty{3}{m}$ だった）
\end{reidai}

\begin{jibun}
速度が $v=t^2+1$（m/s）で、$t=3$ のとき位置が\qty{10}{m}の物体の位置を求めましょう。
\bkjans{$x=\dfrac{t^3}{3}+t+C$、$9+3+C=10$ より $C=-2$。$x=\dfrac{t^3}{3}+t-2$}
\end{jibun}

\begin{reidai}[加速度から位置]
加速度 \qty{-10}{m/s^2}（上向きが＋）で、はじめの位置0、はじめの速度\qty{20}{m/s}の物体の、速度と位置を求めましょう。\par
\kangaekata ③と同じく2回積分し、$C_1$ をはじめの速度で、$C_2$ をはじめの位置で決めます。\par
\shiki $\displaystyle v=\int(-10)\,\dd t=-10t+C_1$。$t=0$ で $v=20$ より $C_1=20$\par
\hspace*{2.6\zw}$\displaystyle x=\int(20-10t)\,\dd t=20t-10\times\frac{t^2}{2}+C_2=20t-5t^2+C_2$\par
\hspace*{2.6\zw}$t=0$ で $x=0$ より $C_2=0$\par
\kotae $v=20-10t$、$x=20t-5t^2$（5-5の投げ上げと同じ式）
\end{reidai}

\begin{jibun}
加速度 \qty{2}{m/s^2}、はじめの速度\qty{3}{m/s}、はじめの位置\qty{1}{m}の物体の、速度と位置を求めましょう。
\bkjans{$v=3+2t$、$x=1+3t+t^2$}
\end{jibun}

\begin{reidai}[いちばん高くなる所]
\bkreidairef{3}のボール（$v=20-10t$、$x=20t-5t^2$）が、いちばん高くなる時刻と、そのときの高さを求めましょう。\par
\kangaekata 上がっているあいだは速度が＋、落ちてくるあいだは $-$ です。＋から $-$ に変わるとき、速度は0になります。その瞬間が、いちばん高い所です（5-7）。\par
\shiki $v=0$ になる時刻\quad $20-10t=0$ より $10t=20$、$t=2$\par
\hspace*{2.6\zw}そのときの高さ\quad $x=20\times2-5\times2^2=40-20=20$\par
\kotae 2秒後、高さ\qty{20}{m}
\end{reidai}

\begin{jibun}
手の高さ\qty{1.5}{m}から、上向きに\qty{10}{m/s}でボールを投げます（$a=-10$）。速度と位置を $t$ の式で表し、いちばん高くなる時刻と高さを求めましょう。
\bkjans{$v=10-10t$、$x=1.5+10t-5t^2$。\mbox{1秒後}、\qty{6.5}{m}}
\end{jibun}

\begin{bkfigure}
\begin{tikzpicture}
\begin{axis}[bkplot, name=g1, width=58mm, height=42mm, xmin=0, xmax=4.5, ymin=0, ymax=25, xtick={1,2,3,4}, ytick={10,20}, xlabel={$t$（s）}, ylabel={$x$（m）}, title={\sffamily\small 位置（高さ）}]
  \addplot[bkcurve, domain=0:4] {20*x-5*x^2};
  \draw[bkGray, dashed, line width=0.5pt] (axis cs:2,0) -- (axis cs:2,20);
  \addplot[only marks, mark=*, mark size=1.8pt, bkInk] coordinates {(2,20)};
\end{axis}
\begin{axis}[bkplot, at={(g1.east)}, anchor=west, xshift=20mm, width=58mm, height=42mm, xmin=0, xmax=4.5, ymin=-24, ymax=24, xtick={1,2,3,4}, ytick={20}, xlabel={$t$（s）}, ylabel={$v$（m/s）}, title={\sffamily\small 速度}]
  \addplot[bkcurve2, domain=0:4] {20-10*x};
  \draw[bkGray, dashed, line width=0.5pt] (axis cs:2,-24) -- (axis cs:2,0);
  \addplot[only marks, mark=*, mark size=1.8pt, bkInk] coordinates {(2,0)};
\end{axis}
\end{tikzpicture}
\bkfigcaption{\bkreidairef{4}のボール。速度が＋から $-$ に変わる2秒（速度0）で、高さがいちばん大きい。}
\end{bkfigure}

\begin{reidai}[ビルから落とす]
高さ\qty{44.1}{m}から、ボールを静かに落とします。地面に着くまでの時間と、そのときの速さを求めましょう（$g=9.8$）。\par
\kangaekata 上向きを＋にすると、④の式で $x_0=44.1$、$v_0=0$ です。$x=44.1-4.9t^2$、$v=-9.8t$。地面は $x=0$。\par
\shiki $44.1-4.9t^2=0$ より $4.9t^2=44.1$、$t^2=44.1\div4.9=9$\par
\hspace*{2.6\zw}2回かけて9になる正の数は3なので $t=3$\par
\hspace*{2.6\zw}$v=-9.8t=-9.8\times3=-29.4$（上向きが＋なので、$-$ は下向き）\par
\kotae 3秒後、下向きに\qty{29.4}{m/s}
\end{reidai}

\begin{jibun}
高さ\qty{19.6}{m}から静かに落としたボールが、地面に着くまでの時間と速さを求めましょう。
\bkjans{$4.9t^2=19.6$ より $t^2=4$、\mbox{2秒後}。$v=-9.8\times2$ で、下向きに\qty{19.6}{m/s}}
\end{jibun}

\noindent\begin{minipage}[c]{0.48\linewidth}
\begin{tikzpicture}
\begin{axis}[bkplot, width=58mm, height=46mm, xmin=0, xmax=3.5, ymin=0, ymax=50, xtick={1,2,3}, ytick={24.5,44.1}, xlabel={$t$（s）}, ylabel={$x$（m）}]
  \addplot[bkcurve, domain=0:3] {44.1-4.9*x^2};
  \addplot[only marks, mark=*, mark size=1.8pt, bkInk] coordinates {(0,44.1) (1,39.2) (2,24.5) (3,0)};
\end{axis}
\end{tikzpicture}
\end{minipage}\hfill
\begin{minipage}[c]{0.5\linewidth}\raggedright
\bkreidairef{5}のボールの高さです。1秒ごとの高さは44.1、39.2、24.5、\qty{0}{m}。1秒ごとに落ちる距離は4.9、14.7、\qty{24.5}{m}で、だんだん大きくなります。速度が1秒ごとに\qty{9.8}{m/s}ずつ速くなるからです。
\end{minipage}

\begin{reidai}[2つの物体が出会う時刻]
物体Aは位置0から、速度 $v=2t$ で＋の向きに動きます。物体Bは位置\qty{15}{m}から、速度\qty{-2}{m/s}で動きます。2つが出会う時刻を求めましょう。\par
\kangaekata 積分と初期条件で、それぞれの位置の式をつくります。出会うのは、位置が等しくなるときです。\par
\shiki A：$x=t^2+C$、$t=0$ で0なので $x=t^2$\par
\hspace*{2.6\zw}B：$x=-2t+C$、$t=0$ で15なので $x=15-2t$\par
\hspace*{2.6\zw}$t^2=15-2t$ より $t^2+2t-15=0$、$(t+5)(t-3)=0$（3-10）\par
\hspace*{2.6\zw}$t=-5$ または $t=3$。時刻は0より後なので $t=3$。位置は $3^2=9$\par
\kotae 3秒後（位置\qty{9}{m}）
\end{reidai}

\begin{jibun}
Aは位置0から速度 $v=4t$ で、Bは位置\qty{30}{m}から速度\qty{-4}{m/s}で動きます。出会う時刻を求めましょう。
\bkjans{A：$x=2t^2$、B：$x=30-4t$。$2t^2=30-4t$、$t^2+2t-15=0$ より\mbox{3秒後}（位置\qty{18}{m}）}
\end{jibun}

\noindent\begin{minipage}[c]{0.48\linewidth}
\begin{tikzpicture}
\begin{axis}[bkplot, width=58mm, height=44mm, xmin=0, xmax=4.2, ymin=0, ymax=18, xtick={1,2,3,4}, ytick={9,15}, xlabel={$t$（s）}, ylabel={$x$（m）}]
  \addplot[bkcurve, domain=0:3.9] {x^2};
  \addplot[bkcurve2, domain=0:3.9] {15-2*x};
  \draw[bkGray, dashed, line width=0.5pt] (axis cs:3,0) -- (axis cs:3,9);
  \addplot[only marks, mark=*, mark size=1.8pt, bkInk] coordinates {(3,9)};
  \node[font=\sffamily\scriptsize, text=bkBlue, anchor=west] at (axis cs:3.92,15.2) {A};
  \node[font=\sffamily\scriptsize, text=bkOrange!80!black, anchor=west] at (axis cs:3.92,7.2) {B};
\end{axis}
\end{tikzpicture}
\end{minipage}\hfill
\begin{minipage}[c]{0.5\linewidth}\raggedright
\bkreidairef{6}の2つの物体の $x$-$t$ グラフです。2本のグラフが交わる点が、同じ時刻に同じ位置にいる「出会い」です。交わる点 $(3,\,9)$ が、計算の答え（3秒後、位置\qty{9}{m}）と合っています。
\end{minipage}

\begin{reidai}[整数をためして解く]
物体Aは位置0から速度 $v=3t^2$ で、物体Bは位置\qty{30}{m}から速度\qty{-1}{m/s}で動きます。2つが出会う時刻を求めましょう（出会う時刻は整数秒です）。\par
\kangaekata 位置の式をつくると、$t^3$ の入った式になり、因数分解の公式が使えません。「整数秒」とわかっているので、$t=1,\,2,\,3,\,\ldots$ を順に入れて、位置が等しくなる時刻をさがします。\par
\shiki A：$x=t^3$（$t^3$ を微分すると $3t^2$）、B：$x=30-t$\par
\hspace*{2.6\zw}$t=1$：Aは1、Bは29\quad $t=2$：Aは8、Bは28\par
\hspace*{2.6\zw}$t=3$：Aは $3^3=27$、Bは $30-3=27$\quad 等しい\par
\kotae 3秒後（位置\qty{27}{m}）
\end{reidai}

\begin{jibun}
Aは位置0から速度 $v=3t^2$ で、Bは位置\qty{12}{m}から速度\qty{-2}{m/s}で動きます。出会う時刻を求めましょう（整数秒です）。
\bkjans{A：$x=t^3$、B：$x=12-2t$。$t=2$ で $8=8$。\mbox{2秒後}（位置\qty{8}{m}）}
\end{jibun}

\noindent\begin{minipage}[c]{0.48\linewidth}
\begin{tikzpicture}
\begin{axis}[bkplot, width=58mm, height=44mm, xmin=0, xmax=4.2, ymin=0, ymax=40, xtick={1,2,3,4}, ytick={10,27}, xlabel={$t$（s）}, ylabel={$x$（m）}]
  \addplot[bkcurve, domain=0:3.35] {x^3};
  \addplot[bkcurve2, domain=0:3.9] {30-x};
  \draw[bkGray, dashed, line width=0.5pt] (axis cs:3,0) -- (axis cs:3,27);
  \addplot[only marks, mark=*, mark size=1.8pt, bkInk] coordinates {(3,27)};
  \node[font=\sffamily\scriptsize, text=bkBlue, anchor=west] at (axis cs:3.38,37.5) {A};
  \node[font=\sffamily\scriptsize, text=bkOrange!80!black, anchor=south] at (axis cs:3.85,28.6) {B};
\end{axis}
\end{tikzpicture}
\end{minipage}\hfill
\begin{minipage}[c]{0.5\linewidth}\raggedright
\bkreidairef{7}のグラフです。Aははじめゆっくりで、だんだん速くなります（速度 $3t^2$）。Bは一定の速さでもどってきます。$t=1,\,2$ ではBのほうが上（先）にいて、$t=3$ で追いつきます。表やグラフで「ためす」のも、りっぱな解き方です。
\end{minipage}

\begin{reidai}[加速度が変わるとき]
加速度が $a=6t$（\unit{m/s^2}）で、はじめの速度\qty{2}{m/s}、はじめの位置0の物体の、速度と位置を求めましょう。2秒後の値も答えましょう。\par
\kangaekata $a$ が変わるので、④の $x=x_0+v_0t+\frac{1}{2}at^2$ は使えません（あれは $a$ が一定のときの式）。③と同じく2回積分します。\par
\shiki $\displaystyle v=\int6t\,\dd t=3t^2+C_1$。$t=0$ で $v=2$ より $C_1=2$\par
\hspace*{2.6\zw}$\displaystyle x=\int(3t^2+2)\,\dd t=t^3+2t+C_2$。$t=0$ で $x=0$ より $C_2=0$\par
\hspace*{2.6\zw}$t=2$ では $v=3\times2^2+2=14$、$x=2^3+2\times2=12$\par
\kotae $v=3t^2+2$、$x=t^3+2t$。2秒後は\qty{14}{m/s}、\qty{12}{m}
\end{reidai}

\begin{jibun}
加速度が $a=12t$（\unit{m/s^2}）で、はじめの速度0、はじめの位置\qty{1}{m}の物体の、速度と位置を求めましょう。1秒後の値も答えましょう。
\bkjans{$v=6t^2$、$x=2t^3+1$。\mbox{1秒後}は\qty{6}{m/s}、\qty{3}{m}}
\end{jibun}

\begin{reidai}[水そうの水位]
水そうの水位が、$\dfrac{\dd h}{\dd t}=0.5-0.1t$（cm/min）の速さで変わります。はじめの水位が\qty{10}{cm}のとき、10分後の水位を求めましょう。\par
\kangaekata いまの水位＝はじめの水位＋変わる速さの積分（⑤）。\par
\shiki $\displaystyle h=10+\int_0^{10}(0.5-0.1t)\,\dd t=10+\Bigl[\,0.5t-0.05t^2\,\Bigr]_0^{10}$\par
\hspace*{2.6\zw}$=10+(0.5\times10-0.05\times10^2)=10+(5-5)=10$\par
\kotae \qty{10}{cm}（はじめの5分は上がり、5分後からは下がって、もとにもどった）
\end{reidai}

\begin{jibun}
お湯の温度が、$\dfrac{\dd T}{\dd t}=3-0.1t$（℃/min）で変わります。はじめが\qty{20}{\degreeCelsius}のとき、10分後の温度を求めましょう。
\bkjans{$20+\Bigl[\,3t-0.05t^2\,\Bigr]_0^{10}=20+25=\qty{45}{\degreeCelsius}$}
\end{jibun}

\begin{reidai}[たまる量がいちばん多くなるとき]
タンクの水の量が $\dfrac{\dd V}{\dd t}=30-6t$（L/min）で変わり、はじめは\qty{50}{L}です。水の量がいちばん多くなるのは何分後で、そのとき何Lですか。\par
\kangaekata 変わる速さが＋のあいだはふえ、$-$ になるとへります。速さが0になる所で、ふえるからへるに変わります（5-7、\bkreidairef{4}と同じ考え）。\par
\shiki $\displaystyle V=\int(30-6t)\,\dd t=30t-3t^2+C$。$t=0$ で $V=50$ より $C=50$\par
\hspace*{2.6\zw}変わる速さ $30-6t=0$ より $t=5$\par
\hspace*{2.6\zw}$V=50+30\times5-3\times5^2=50+150-75=125$\par
\kotae 5分後、\qty{125}{L}
\end{reidai}

\begin{jibun}
水位の変わる速さが $4-2t$（cm/min）で、はじめの水位は\qty{10}{cm}です。水位がいちばん高くなる時刻と水位を求めましょう。
\bkjans{$h=10+4t-t^2$。$4-2t=0$ より\mbox{2分後}、\qty{14}{cm}}
\end{jibun}

\begin{reidai}[あとの時刻の様子から、さかのぼる]
ボールを真上に投げました（$a=-10$、上向きが＋）。投げてから2秒後に、いちばん高い\qty{25}{m}の所に来ました。はじめの速度 $v_0$ と、投げた高さ $x_0$ を求めましょう。\par
\kangaekata 初期条件が「2秒後」の値で2つ（いちばん高い所なので速度0、高さ25）わかっています。②のように、わかっている時刻の値を入れて $C_1$、$C_2$ を決めます。$t=0$ を入れれば、$v_0$ と $x_0$ がわかります（$C_1=v_0$、$C_2=x_0$）。\par
\shiki $\displaystyle v=\int(-10)\,\dd t=-10t+C_1$\par
\hspace*{2.6\zw}$t=2$ で $v=0$ なので $0=-10\times2+C_1$、$C_1=20$\par
\hspace*{2.6\zw}$\displaystyle x=\int(-10t+20)\,\dd t=-5t^2+20t+C_2$\par
\hspace*{2.6\zw}$t=2$ で $x=25$ なので $25=-5\times4+20\times2+C_2$\par
\hspace*{2.6\zw}$25=20+C_2$、両辺から20をひいて $C_2=5$\par
\hspace*{2.6\zw}確かめ：$x=5+20t-5t^2$ に $t=2$ を入れると $5+40-20=25$\par
\kotae はじめの速度\qty{20}{m/s}（上向き）、投げた高さ\qty{5}{m}
\end{reidai}

\begin{jibun}
ボールを真上に投げると（$a=-10$）、1秒後に、いちばん高い\qty{6.5}{m}の所に来ました。はじめの速度と、投げた高さを求めましょう。
\bkjans{$v=-10t+C_1$、$0=-10+C_1$ より $C_1=10$。$x=-5t^2+10t+C_2$、$6.5=-5+10+C_2$ より $C_2=1.5$。はじめの速度\qty{10}{m/s}、高さ\qty{1.5}{m}}
\end{jibun}

%% ---------------- まちがえやすいところ・物理工学 ----------------
\bksection{tsumazuki}{まちがえやすいところをチェック}

\begin{tsumazuki}{$C$ を、いつも0にしてしまう}
\bkng{$v=2t$、$t=0$ で位置\qty{3}{m} → $x=t^2$}\bkok{$x=t^2+C$、$C=3$ → $x=t^2+3$}
$x=t^2$ とすると、はじめの位置を0と決めつけたことになります。初期条件がわかればそれで $C$ を決め、なければ $+C$ を残します。
\end{tsumazuki}

\begin{tsumazuki}{2回積分するのに、$C$ を1つしか使わない}
加速度から位置までは、$C_1$ と $C_2$ の2つが出ます。$C_1$ ははじめの速度、$C_2$ ははじめの位置で決めます。1つの $C$ にまとめてはいけません。
\end{tsumazuki}

\begin{tsumazuki}{速度が変わるのに「位置＝速度×時間」とする}
\bkng{$v=2t$ の3秒後は $v=6$ なので $x=6\times3=\qty{18}{m}$}\bkok{$x=\displaystyle\int_0^3 2t\,\dd t=\qty{9}{m}$}
「速度×時間」は、速度が一定のときだけ使えます。3秒後の速さ6でずっと走っていたわけではないからです。
\end{tsumazuki}

\begin{butsuri}
\begin{itemize}[leftmargin=1.2\zw, itemsep=1pt, topsep=0pt, label={\color{bkPurple}●}]
\item \textbf{慣性航法}\\飛行機やロケットは、加速度センサーの値を2回積分して位置を計算します。はじめの位置と速度を正確に入れておきます。
\item \textbf{物理の法則と初期条件}\\物理の法則は「変化の速さ」（運動なら加速度）を決めます。法則とはじめの状態がそろって、はじめて先のようすが1つに決まります。これが第6巻の微分方程式の考え方です（6-1）。
\item \textbf{ダム}\\入る水と出る水の差を積分して、はじめの水位にたし、これからの水位を予想します。
\end{itemize}
\end{butsuri}

%% ---------------- 章末問題 ----------------
\bkendproblems
\bkproblemgroup{基本}{3}

\bkproblem 速度が $v=6t$（m/s）で、$t=0$ の位置が\qty{4}{m}の物体の位置を求めましょう。
\bkpans{$x=3t^2+4$}

\bkproblem 速度が $v=4t+1$（m/s）で、$t=0$ の位置が\qty{-3}{m}の物体の位置を求めましょう。
\bkpans{$x=2t^2+t-3$}

\bkproblem 速度が $v=3t^2-2t$（m/s）で、$t=2$ のとき位置が\qty{5}{m}の物体の位置を求めましょう。
\bkpans{$x=t^3-t^2+C$、$8-4+C=5$ より $C=1$。$x=t^3-t^2+1$}

\bkproblem 高さ\qty{2}{m}から、上向きに\qty{15}{m/s}でボールを投げます（$a=-10$）。速度と位置を $t$ の式で表し、いちばん高くなる時刻と高さを求めましょう。
\bkpans{$v=15-10t$、$x=2+15t-5t^2$。\mbox{1.5秒後}、\qty{13.25}{m}}

\bkproblem 高さ\qty{78.4}{m}から静かに落としたボールが、地面に着くまでの時間を求めましょう（$g=9.8$）。
\bkpans{$4.9t^2=78.4$ より $t^2=16$、\mbox{4秒後}}

\bkproblem 高さ\qty{4.9}{m}から静かに落としたボールが、地面に着くまでの時間と、そのときの速さを求めましょう（$g=9.8$）。
\bkpans{$4.9t^2=4.9$ より $t^2=1$、\mbox{1秒後}。下向きに\qty{9.8}{m/s}}

\bkproblem 加速度\qty{3}{m/s^2}、はじめの速度\qty{4}{m/s}、はじめの位置0の物体の、2秒後の速度と位置を求めましょう。
\bkpans{$v=4+3t$、$x=4t+1.5t^2$。\qty{10}{m/s}、\qty{14}{m}}

\bkproblem お湯の温度が $\dfrac{\dd T}{\dd t}=2-0.2t$（℃/min）で変わります。はじめが25\,℃のとき、10分後の温度を求めましょう。
\bkpans{$25+\Bigl[\,2t-0.1t^2\,\Bigr]_0^{10}=25+10=\mbox{$35$\,℃}$}

\bkproblemgroup{チャレンジ}{3}

\bkproblem 加速度が $a=6t$（\unit{m/s^2}）で、はじめの速度\qty{-3}{m/s}、はじめの位置0の物体の、速度と位置を求めましょう。2秒後の値も答えましょう。
\bkpans{$v=3t^2-3$、$x=t^3-3t$。2秒後は\qty{9}{m/s}、\qty{2}{m}}

\bkproblem ロケットの加速度が $a=2t$（m/s$^2$）でふえていきます。はじめの速度と位置は0です。6秒後の速度と位置を求めましょう。
\bkpans{$v=t^2$、$x=\dfrac{t^3}{3}$。\qty{36}{m/s}、\qty{72}{m}}

\bkproblem タンクの水の量が、$\dfrac{\dd V}{\dd t}=20-4t$（L/min）で変わり、はじめは\qty{100}{L}です。水の量がいちばん多くなるのは何分後で、そのとき何 L ですか（5-7）。
\bkpans{$V=100+20t-2t^2$。$20-4t=0$ より\mbox{5分後}、\qty{150}{L}}

\bkproblem 水位が $\dfrac{\dd h}{\dd t}=6-2t$（cm/min）で変わり、はじめの水位は\qty{5}{cm}です。5分後の水位と、水位がいちばん高くなる時刻と水位を求めましょう。
\bkpans{$h=5+6t-t^2$。5分後は\qty{10}{cm}。$6-2t=0$ より\mbox{3分後}、\qty{14}{cm}}

\bkproblem 物体Aは位置0から速度 $v=2t$ で、物体Bは位置\qty{8}{m}から速度\qty{2}{m/s}で、同じ向きに動きます。AがBに追いつく時刻と位置を求めましょう。
\bkpans{A：$x=t^2$、B：$x=8+2t$。$t^2-2t-8=0$、$(t-4)(t+2)=0$ より\mbox{4秒後}（位置\qty{16}{m}）}

\bkproblem 物体Aは位置0から速度 $v=3t^2$ で、物体Bは位置\qty{16}{m}から速度\qty{-4}{m/s}で動きます。2つが出会う時刻を求めましょう（出会う時刻は整数秒です）。
\bkpans{A：$x=t^3$、B：$x=16-4t$。$t=2$ でAは8、Bは $16-8=8$。\mbox{2秒後}（位置\qty{8}{m}）}

\bkproblem ボールを真上に投げると（$a=-10$、上向きが＋）、3秒後に、いちばん高い\qty{50}{m}の所に来ました。はじめの速度と、投げた高さを求めましょう。
\bkpans{$v=-10t+C_1$、$0=-30+C_1$ より $C_1=30$。$x=-5t^2+30t+C_2$、$50=-45+90+C_2$ より $C_2=5$。はじめの速度\qty{30}{m/s}、高さ\qty{5}{m}}

\bkproblem 速度が $v=2t+1$（m/s）の物体が、$t=2$ のとき位置\qty{10}{m}にいました。$t=0$ の位置と、$t=4$ の位置を求めましょう。
\bkpans{$x=t^2+t+C$、$4+2+C=10$ より $C=4$。$t=0$ で\qty{4}{m}、$t=4$ で $16+4+4=\qty{24}{m}$}

\begin{furikaeri}
\begin{checklist}
\item 初期条件から、積分定数を決められる。\hfill{\footnotesize\sffamily\color{bkTheme}→\pageref{bk:5-13:shoki}ページ}
\item 加速度から、速度と位置を求められる。\hfill{\footnotesize\sffamily\color{bkTheme}→\pageref{bk:5-13:toukasoku}ページ}
\item 「いまの量＝はじめの量＋変化の積分」を使える。\hfill{\footnotesize\sffamily\color{bkTheme}→\pageref{bk:5-13:teisekibun}ページ}
\end{checklist}
\end{furikaeri}

\bkfukidashi{delta}{第2部はここまで。次の部では、$e$ や $\sin$・$\cos$ という特別な関数を、微分・積分してみるよ。}
